\documentclass[twocolumn]{aastex631}
\begin{document}
\title{The reionization ionizing-photon-budget shortfall for star-forming galaxies is not robust to systematics at z\~{}8 (jwsthigh systematic corner)}
\author{NebulaMind Lab (autonomous pipeline)}
\affiliation{NebulaMind Lab, \url{https://lab.nebulamind.net}}
\begin{abstract}
Reionization ionizing-photon-budget reconciliation at z\~{}8: star-forming galaxies require f\_esc=0.133 (+0.126/-0.064) to close the budget (Madau-Dickinson SFRD, log xi\_ion=25.5+/-0.15, clumping C=2-5, JWST-SFRD tail), versus indirect-proxy-inferred f\_esc=0.062 (+0.110/-0.039) from LzLCS O32/beta calibrations. Median delta(required-inferred)=+0.061 dex-frac (16-84\%: -0.049 to +0.189); 74\% of the systematic MC shows a shortfall. Conclusion: the budget CLOSES within the systematic, and the sign holds under both O32 and beta calibrations. The result is bounded by the xi\_ion x clumping x proxy-calibration systematic, not by statistics. Generated autonomously from public data (jwst) via the NebulaMind Lab runner: a bounded, reproducible, descriptive study.
\end{abstract}
\section{Introduction}
Recent studies have highlighted a potential crisis in understanding the reionization process, with concerns about whether star-forming galaxies can produce enough ionizing photons to account for the observed reionization [Muñoz2024]. This has led to increased scrutiny of the ionizing photon budget and the factors that influence it. Previous work has explored various aspects of this problem, including the role of galaxy ionizing photon budgets at high redshifts [Duncan2015] and the challenges posed by absorption-dominated reionization scenarios [Davies2021]. To address these questions, researchers have employed a range of models and approaches, such as excursion set reionization models [Park2022] and analytic methods [Madau2017].
\section{Data and method}
In this study, we adopt a literature-anchored budget calculation approach to investigate the ionizing photon budget at z\~{}8. We use the Madau \& Dickinson (2014) analytic fitting function for the cosmic star formation rate density (SFRD). For the ionization parameters xi\_ion and O32/beta f\_esc proxy calibrations, we rely on published values from LzLCS [Chisholm+22, Flury+22; Simmonds+24]. Importantly, our analysis does not utilize survey catalog data or observational data from JWST, SDSS, or TNG. Instead, it focuses on reconciling systematic uncertainties in the literature to determine whether star-forming galaxies can account for reionization.
\section{Result}
Our calculations reveal that to close the ionizing photon budget at z\~{}8, star-forming galaxies must have an escape fraction f\_esc of 0.133 (+0.126/-0.064). This value is higher than the indirect-proxy-inferred f\_esc of 0.062 (+0.110/-0.039) derived from LzLCS O32/beta calibrations. The median difference between the required and inferred escape fractions is +0.061 dex-frac, with a range of -0.049 to +0.189 (16-84\% confidence interval). Notably, 74\% of our systematic Monte Carlo simulations indicate a shortfall in ionizing photons.
\begin{figure}\centering\includegraphics[width=\columnwidth]{result.png}\caption{The reionization ionizing-photon-budget shortfall for star-forming galaxies is not robust to systematics at z\~{}8 (jwsthigh systematic corner)}\end{figure}
\section{Caveats}
It is essential to acknowledge the limitations of our approach. Our analysis relies on automated, single-selection, uncalibrated measurements, which may introduce biases and uncertainties. The results are sensitive to the choice of xi\_ion, clumping factor C, and proxy calibration, highlighting the need for further research to refine these parameters. Additionally, our study does not account for potential variations in the SFRD or other factors that could influence the ionizing photon budget. These caveats emphasize the importance of continued investigation into the complex processes driving reionization.
\section*{References}
\footnotesize
[Muoz2024] Muñoz et al. (2024). Reionization after JWST: a photon budget crisis?. bibcode 2024MNRAS.535L..37M \par [Park2022] Park et al. (2022). Calibrating excursion set reionization models to approximately conserve ionizing photons. bibcode 2022MNRAS.517..192P \par [Duncan2015] Duncan et al. (2015). Powering reionization: assessing the galaxy ionizing photon budget at z < 10. bibcode 2015MNRAS.451.2030D \par [Davies2021] Davies et al. (2021). The Predicament of Absorption-dominated Reionization: Increased Demands on Ionizing Sources. bibcode 2021ApJ...918L..35D \par [Madau2017] Madau (2017). Cosmic Reionization after Planck and before JWST: An Analytic Approach. bibcode 2017ApJ...851...50M

\end{document}
